\documentclass{article}
\usepackage[T1]{fontenc}
\usepackage{amsmath,amssymb,amsthm}
\usepackage{graphicx,booktabs,microtype,xurl}
\usepackage{mlsys2025}
\usepackage[hidelinks]{hyperref}
\urlstyle{same}
\hypersetup{pdfsubject={Working research draft, September 2026}}
\makeatletter
\renewcommand{\Notice@String}{Working draft, 28 September 2026. Not submitted.}
\makeatother
\newtheorem{proposition}{Proposition}
\newtheorem{lemma}[proposition]{Lemma}
\mlsystitlerunning{Private Model Export for New Tasks}
\begin{document}
\twocolumn[
\mlsystitle{Private Model Export for New Tasks}
\begin{mlsysauthorlist}
\mlsysauthor{Anonymous Authors}{anon}
\end{mlsysauthorlist}
\mlsysaffiliation{anon}{Anonymous Institution}
\mlsyscorrespondingauthor{Anonymous Author}{anonymous@example.invalid}
\vskip 0.3in
\begin{abstract}
An exported model cannot be withdrawn from recipients who retain or share
it. We study how to release a model for another task while preserving
the earlier output and bounding the privacy of their complete history.
For a known finite old channel, we optimize the new release subject to
its fixed old marginal. An exact construction handles independent uniform
binary tasks under all-pairs input adjacency; with old binary randomized
response it takes $O(r^2)$ operations for $r$ new coordinates. Our central
record-level construction instead uses a finite linear program. Count
and checked label symmetries reduce this program, and rational certificates
verify privacy and bound numerical suboptimality.

Joint optimization improves on an old-output-conditioned composition
baseline in twelve finite synthetic configurations. Checked symmetry
certificates avoid expanding the proof; fixed disjoint blocks extend
execution to ensembles without enlarging each optimization problem.
A central protocol binds the actual channels to one continuation from a
menu frozen before protected input, on unchanged data. On 75 real-data
cohorts from three corpora, however, the finite optimizer matches reuse
and trails the stronger public learner by 0.02597 mean accuracy.
A public information-value bound explains this failure and rules out
the prescribed gain before solving the private optimization problem.
The result is a certified finite continuation method and a test for
when it is unnecessary, not demonstrated practical gains for general
model training or an agency deployment.
\end{abstract}
]
\printAffiliationsAndNotice{}
\hypersetup{pdftitle={Private Model Export for New Tasks}}
\raggedbottom
\clubpenalty=10000
\widowpenalty=10000
\displaywidowpenalty=10000

\section{Introduction}
A central public agency may oversee several generations of models from citizen
records whose sensitivity outlasts the original service. Recipients
can retain exported weights, share them with other recipients and make
local predictions. The default observer therefore pools every model
released by the covered process, across tasks, datasets and recipients.
A new task changes what the agency must build without removing what
this group already knows. Removing direct
identifiers does not prevent inference from linked releases
\citep{ganta2008composition}; replacing an agency's current model does
not erase an earlier exported copy.

The technical problem is the following: \emph{given the conditional law
of what has already been exported, what is the best achievable new-task
accuracy at a fixed bound on the complete release history?} A new mechanism
must preserve the actual earlier output and its law. A scalar privacy
charge alone can lose information relevant to this decision: different
old channels with the same bound can leave different room for a new task.
The challenge is to characterize that room and construct an attaining
extension under explicit assumptions. This paper studies finite model
families for which that law is available, not an analyzer for arbitrary
trained weights. In the implemented protocol, ``new task'' means a later
release from a predeclared menu: its protected labels already exist when
the first model is sampled.

The old disclosure may be privatized training records or only an
internally trained model. Our analysis includes central computation for
the new task while preserving either kind of earlier disclosure.
Singapore GovTech documents synthetic-data sharing for AI development;
the U.S. Census SIPP report documents requests for new variables and
validation after synthetic publication
\citep{govtech2026syntheticapplications,benedetto2013ssb}.
These establish related operational needs, not deployment of our
finite-channel extension. Its agency use remains proposed.

\input{related-work-revised.tex}
\subsection{Our contribution}
\label{sec:contribution}
We make three contributions within this finite-law setting.
First, we derive a finite-breakpoint extension algorithm for independent
uniform new binary tasks and a known finite old release channel. The
binary-RR instance takes $O(r^2)$ arithmetic operations for $r$ new
coordinates. Its single-new-coordinate case is prior work; the extension
concerns multiple new coordinates and general finite old channels.

Second, we implement count and checked class-label reductions for known
finite old laws, including multiclass outputs and binary-feature
classifiers. Compact rational certificates verify privacy and bound
suboptimality without expanding the probability table. A stronger
composition control conditions on the actual old output. Joint design
improves on it in all twelve tested synthetic configurations; a real-data
evaluation across 75 cohorts instead finds no gain over reuse. We explain
that shortfall using a public value-of-information bound, an elementary
decision-theoretic check rather than a new DP theorem.

Third, the release pipeline binds that public specification and the actual
channel menu before protected input. It reserves the complete history
bound and delivers one immutable continuation. Fixed disjoint blocks
reuse a certified template to construct ensembles under standard parallel
composition. Reusing identical public checks reduces repeated work while
current authorization remains uncached. We measure these costs separately
from utility: scalable execution cannot make an uninformative model useful.
The supplement retains the essential proofs and methods; complete results
remain in the experimental companions.

\input{preliminaries.tex}
\section{Our Release Protocol}
\label{sec:protocol-evaluation}
Our protocol assumes a central public agency that maintains one
authoritative dependency graph for the entire dataset under its
responsibility. The graph covers the original data, derived dataset
versions, training computations and model releases. This setting
corresponds to an agency allowing public access to approved versions
of its data for model training. The agency audits proposed releases
and maintains a common privacy account across users and projects.

At ingestion, direct identifiers are removed and a separate restricted
linkage boundary supplies stable pseudonyms. Pseudonymization does not
make records anonymous or prevent inference from other fields. External
users receive only accounted private releases; both data and models enter
the pooled history. The ingestion and vault requirements remain specified
architecture, not deployment evidence (Appendix~\ref{app:stage-zero}).

\begin{figure*}[!t]
\centering
\includegraphics[width=\textwidth]{figures/central-audit-graph.pdf}
\caption{Agency dataset and release history. Both ingestion paths
remove direct identifiers before analysis; a separate vault controls
identity lookup, without preventing inference from other fields.
Registering $B$ creates a
new main-dataset version; every derived dataset links to its source
version. Models $R_1,R_2$ reuse $A_1$, while $R_4$ combines $A_2,A_3$.
Solid arrows show dependencies; the dashed line shows shared people.
Prior disclosures and charges persist. New data require reviewed
membership, contribution bounds and privacy evidence. Illustrated
external training uses only accounted DP releases and fixed context.
The graph can register changing datasets, but the finite continuation
analyzed here uses one unchanged data state. The vault and ingestion
controls are specified, not deployment evidence.}
\label{fig:release-protocol}
\end{figure*}

\subsection{Register data, computations and every release}
\label{sec:central-graph}
Write the agency's directed dependency graph as
$G_t=(\mathcal D_t\cup\mathcal C_t\cup\mathcal R_t,\mathcal E_t)$.
Dataset nodes $\mathcal D_t$ identify immutable versions and their
protected units. Each derived dataset links to the relevant version of
the agency's main dataset. An addition creates a registered successor
version and updated membership records, while retaining the earlier
version and its release history (Figure~\ref{fig:release-protocol}).
Computation nodes $\mathcal C_t$ identify protected
accesses, transformations, randomness dependencies and privacy evidence.
Release nodes $\mathcal R_t$ identify each shared private dataset,
exported model package or other observable disclosure, including a
refusal without a model, and
its authorization. Edges state which inputs produced each result;
every retraining or deletion update adds a new version. Re-delivery of
a saved package refers to its existing computation, rather than silently
creating another draw under the same identity.

The agency also maintains confidential unit-membership links across
dataset nodes. Separate branches need not contain different people.
Dataset names and recipients do not alter this account. Unresolved
overlap requires a conservative bound or withheld authorization. Pairwise
overlap counts do not establish each person's complete exposure; unit
membership or a sound higher-order bound is required. These confidential
links are fixed in our adjacency and are not public exports.

\subsection{Central audit and authorization}
\label{sec:central-audit}
For a proposed release, the agency traces its ancestors and the prior
release history. It checks three matters before authorization.

\paragraph{Complete scope and credible execution.}
Reconcile declared inputs, protected units and existing releases with
the controlled data and export inventory. Include protected labels,
preprocessing, selection, validation and visible audit outcomes.
Bind the proposed package to the audited data version and computation.
A submitted graph or artifact hash alone does not establish truthful
execution. Appendix~\ref{app:central-audit} separates these obligations
from the local prototype's evidence.

\paragraph{A valid joint privacy argument.}
Audit the whole dependency path, not a model's attack score. Reuse is
free only when outputs are post-processing of an already-accounted
private source. Fresh protected access needs conditional composition
or an explicit joint-channel proof; correlated marginal guarantees
cannot simply be added. These arguments must hold over the declared
neighboring datasets and permitted histories, not only the realized
model bytes. Charge each certified source for each affected
unit once, not once per descendant model. The finite constructions
below provide particular joint arguments. Missing privacy evidence
blocks release; an audit cannot retrospectively repair an earlier
invalid disclosure.
The charge must cover the unit's full contribution to each source.

\paragraph{Task utility.}
Freeze the package, prediction inputs and task requirements before
independent evaluation. Passing a task requirement and passing privacy
checks are separate requirements; a useful model or unsuccessful attack
does not certify DP. Changed data, including deletion, require a new
mechanism argument; registration alone does not make a fixed-data
extension applicable.

\subsection{Extend a known finite release channel}
\label{sec:commitment}
Suppose the pooled earlier output $O$ has known conditional law
$P(o\mid s)$. It includes disclosures to every covered recipient;
the finite construction can treat this bundle as one old channel.
A new workload requires $r$ independent uniform binary coordinates
$Z_1,\ldots,Z_r$, independent of the old protected input $s$.
The agency may access the data again, but must retain the actual old
output. The complete pair of old output and new decisions must satisfy
pure $E$-DP under replacement of $(s,Z)$. We optimize minimum new-task
accuracy, averaging over a declared public prior on $s$.

This construction treats $(s,Z)$ as one protected unit and compares
every pair of its values. Exposing its separately randomized row outputs
gives the local-DP observation model, even when the agency samples them.
It answers how to extend an irrevocable finite disclosure, rather than
which mechanism the agency should choose before any disclosure exists.
For population learning, a central computation that never exposes
row-level responses is a necessary comparator. The model-history
formulation below handles a different, general record-neighbor graph;
its extension principle does not require a local old release.

Put $N=2^r$ and $\lambda=e^E$. Symmetry groups new error vectors by
their number $h$ of incorrect coordinates. For each old output choose
a probability floor
\[
 \frac{\max_sP(o\mid s)}{N\lambda}
 \le m_o\le \frac{\min_sP(o\mid s)}N.
\]
For each old state, assign that floor to every error vector and fill
remaining probability in increasing order of $h$, up to $\lambda m_o$
per vector. The optimum occurs at an endpoint or one of finitely many
scalar breakpoints. Proposition~\ref{prop:finite-commitment} gives the
formula, boundary cases and proof. For binary randomized response (RR)
as the old channel, the calculation takes $O(r^2)$ arithmetic operations.
It requires the complete old channel, not merely its privacy scalar.
It is not a tractable characterization of arbitrary trained weights.

The one-new-coordinate formula recovers prior work
\citep{yamamoto2024privacyoptimized}. For several coordinates the
proper comparison is an optimized joint mechanism for the new vector,
independent of the old response, not only coordinate-wise RR.
\subsection{Use the actual retained model history}
\label{sec:weights-main}
Let $P_D(m)$ be the joint law of the pooled old history, including
weights and other data-dependent disclosures to any covered recipient.
Separate models' marginal laws do not determine this joint law.
An extension $Q_D(m,a)$
must satisfy
\[
 \begin{aligned}
 \sum_a Q_D(m,a)&=P_D(m),\\
 Q_D(m,a)&\le e^E Q_{D'}(m,a)\quad(D\sim D').
 \end{aligned}
\]
Given the actual $m$, the agency can realize a feasible finite extension
by sampling $a$ from $Q_D(m,a)/P_D(m)$ on its positive support. This
requires protected data access and knowledge of the conditional laws;
it is not free post-processing of $m$.

Summing the joint inequality over $a$ shows that the old law must itself
satisfy the requested bound. A continuation cannot repair a nonprivate
earlier release. For example, a deterministic old output that differs on
neighbors already gives an event of probability one versus zero and
cannot admit a finite pure-DP extension. Knowing a law is therefore
necessary but not sufficient for feasibility.

Protecting a full private training representation is sufficient but can
be unnecessarily restrictive when only particular weights were exported
and no recipient obtained that representation.
If private training data releases are shared, they remain part of the old
disclosure history; a model-only comparison cannot omit them.
Appendix~\ref{sec:actual-model-history} proves a strict separation in a
three-record learning problem and a recoverability condition for equality.
An omitted disclosure invalidates the comparison. The result is not a
budget refund or a general model analyzer; ordinary composition remains
a conservative baseline.

For a fixed $q$-value record alphabet, suppose the old channel and the
utility objective are invariant to record permutations. Averaging a
feasible extension over those permutations preserves its old marginal,
record-level privacy and utility. The channel can then depend only on
the $q$ category counts, with $\binom{n+q-1}{q-1}$ states; a neighboring
state moves one record between categories. Proposition~\ref{prop:count-reduction}
proves the reduction and gives the necessary multiplicity weights.
For our two-bit records and two binary model outputs, this gives
$4\binom{n+3}{3}$ variables instead of $4^{n+1}$.
It requires invariance of the complete earlier disclosure, not merely
an informal claim that a model ignores row order. Count aggregation
removes row-order redundancy, but still covers every possible count
state, rather than just the observed count vector. Both the record count
$n$ and alphabet size $q$ limit tractability.

With uncertain public utility laws, optimize a mixture or worst-case
regret over a declared set; the privacy constraints stay unchanged.
Choosing a task from protected data instead needs a separate charge.
Appendix~\ref{app:uncertain-workloads} proves the bound for a fixed menu
of joint channels and a private selector after the old export.



\subsection{Check whether private adaptation can help}
\label{sec:information-value}
Before reserving a continuation, compare its potential utility with
public prediction and reuse. For the declared average linear objective
$U(Q)=c_0+\sum_{D,m,a}c_{Dma}Q_D(m,a)$, define
\begin{align}
 U_{\rm reuse}&=c_0+\sum_m\max_a\sum_D P_D(m)c_{Dma},\nonumber\\
 U_{\rm full}&=c_0+\sum_{D,m}P_D(m)\max_a c_{Dma}.
 \label{eq:information-value}
\end{align}
The first is the best prediction rule using only the old output.
The second allows the rule to see the complete private state without
any privacy restriction. If $U_{\rm full}-U_{\rm reuse}<\Delta$ for a
prespecified required gain $\Delta>0$, no private continuation in this
family can meet that gain. This check needs no LP, although it still
enumerates the finite public problem. Otherwise a certified private
upper bound can sharpen the check. Passing either check only leaves
improvement possible; it does not authorize release.

When the check rules out adaptation, the agency can construct the
old-only rule directly and certify its joint history at the old law's
bound. Choosing it before reservation can avoid an unnecessarily large
charge. This is tighter prospective accounting, not stronger privacy
for an identical channel or a refund after release. The conclusion is
conditional on the public utility law and model class; it says nothing
about arbitrary unknown populations (Appendix~\ref{app:information-value}).

\subsection{Commit the release and audit subsequent changes}
\label{sec:workload-extension}
\label{sec:implemented-assurance}
Reserve the required allowance before the approved protected computation
and commit the audited release record before external access. Concurrent
proposals must fit the common allowance jointly. Only the authorized
package may leave the controlled boundary. Changes to data, code,
private inputs, noise relationships or disclosed metadata require a new
audit. Reconcile actual exports with the graph; preserve prior release
nodes after revocation or deletion. These are central control requirements,
not evidence of deployed multi-agency enforcement.

Our bounded instantiation admits independent private sources, exact reuse
and a checked finite-model continuation on fixed registered memberships.
Before accepting protected inputs, it freezes the complete public problem,
model meanings, selection kernel and the actual certified channel menu.
The optimizer uses only that public specification. Exact checks bind each
candidate's old marginal, adjacency, requested odds and objective. Merely
fixing the optimization problem is insufficient: privately choosing between
two valid solutions could reveal data. The whole history allowance is
reserved before the old draw. After any accounted
task selection, the executor samples the new model conditional on the
saved old output and unchanged confidential counts. One extension slot
and immutable committed models prevent retries or other recipients from
obtaining uncharged draws. At total odds $3$, the conservative rational
charge is $11/10>\log3$; unused or revoked reservations are not refunded.
This covers a prospectively frozen menu on the same data state, not arbitrary later
tasks or changed records. Appendix~\ref{app:finite-enforcement} gives
the execution argument.

For larger rosters, a fixed public partition can assign each privacy unit
to exactly one block, with independent mechanism randomness across
blocks. Each block mechanism uses only that block's protected values
and fixed public context. The executor applies the same certified finite
template to each block and exports their heads as an ensemble. A changed unit affects
only one complete block history; the pooled bound is the largest block
bound, not their sum. Voting is post-processing. All contributions from
one protected unit must remain in that block, and new partitions or
overlapping histories cannot reset its account. This uses standard
parallel composition, not optimal joint learning on the full roster
(Appendix~\ref{app:block-execution}).

Repeated verification of identical public certificates and plan templates
can be cached. Cache identity includes the full canonical inputs,
verification limits, model meanings and current mathematical charge.
Source binding, reservation, revocation and release authority are checked
against current state on every request. A cache hit is never itself an
authorization.

Other correlated-noise declarations remain unsupported.
Appendix~\ref{app:central-gate-proof} proves the account invariant under
truthful mechanisms, complete mediation and durable transitions. Graph
and ledger representations give equivalent admission under the same evidence.

Evaluate frozen candidates with independent public evidence or an
accounted private procedure, including selection. Identify what evidence
could resolve a threshold-crossing uncertainty; failure of tested
candidates is not general infeasibility.



\section{Experiments}
\label{sec:experiments}
The primary experiments ask how much utility is lost by keeping an old
export, how much a joint extension recovers over conditional composition,
and whether the resulting model can be certified and released at reasonable
cost. For a higher-is-better score $U$, we report
\[
 C=U_{\rm no\ old}-U_{\rm joint},\qquad
 G=U_{\rm joint}-U_{\rm conditional}.
\]
These utility comparisons score the new prediction. The no-old arm is
a counterfactual with no previous disclosure and the
same total allowance; it cannot be added to an actual old history for
free. The conditional arm optimizes a new learner separately for every
old output, using that output's worst-case residual privacy allowance;
the independent arm cannot condition on it. We first test finite
families and uncertain workloads, then exact certification and execution,
and finally held-out prediction with the finite optimizer on public
datasets. Latent-task accuracy, prediction accuracy and runtime answer
different questions. Certificates, timings and cohort-level comparisons
provide distinct kinds of evidence.

\paragraph{Finite reference checks.}
The earlier scalar solver checks 81 RR and 108 general-channel cells;
63 small unrestricted LPs agree to $2.9\times10^{-13}$.
The three-record model-only optimum is $0.739665$, versus $0.5$
when the full old randomized cache is exposed and $0.782$ with no
old export. The full numerical records remain in the experimental companions.
The following studies widen the model and workload assumptions.
\input{count-extension-results.tex}
\input{richer-model-results.tex}
\input{uncertain-workload-results.tex}
\input{finite-enforcement-results.tex}
\input{finite-realdata-results.tex}

\section{Discussion and Guarantees}
\label{sec:discussion}
\subsection{What follows from the written arguments}
The shell construction is optimal for its all-pairs input adjacency
(Appendix~\ref{app:joint-commitment}); the count formulation instead
uses record adjacency. Permutation averaging preserves its optimal
utility under the stated invariance conditions. Additional checked
label symmetries preserve that optimum, while rational
certificates verify privacy and bound suboptimality
(Appendices~\ref{app:count-extension} and~\ref{app:compact-certification}).
The new compact certificate checks the symmetry and a quotient dual;
averaging proves its bound for the unrestricted finite problem
(Appendix~\ref{app:compact-certification}). The public state domain is
still explicit, so the count-space growth exponent is unchanged.
The old-model law must be known. Fixed-block parallel composition and
public verification caches improve execution, not that mathematical limit.
The graph alone does not prove privacy.

\subsection{What the experiments do not establish}
The public-law uncertainty study changes utility assumptions while
retaining known mechanism laws; it does not infer arbitrary old-model
laws from data. The real-data experiment now applies the finite optimizer
and the stronger controls. The separate information-value bound rules
out the prescribed gain for these finite public objectives, not useful private adaptation in richer
families. The coarse feature, tiny private sample and near-constant best
action jointly limit this design. Adding more independent copies of that
same head does not repair the loss. A positive practical claim would need
a richer tractable family, action-changing private information and new
evaluation populations. Repeating noise or tuning on these cohorts is
insufficient.

A changed-data extension still needs a retained-data reference,
deletion-respecting computation and a pooled-history bound.

\subsection{Privacy and institutional scope}
\label{sec:privacy-basis}
The active experiments protect records. Person-level deployment requires
complete contribution bounds and truthful linkage. The account covers
all recipients within the participating boundary, not unknown outside
models. Authenticated execution, complete export mediation, vault security
and protection against account forks remain deployment obligations
(Appendix~\ref{app:central-audit}).

The prototype executes a frozen finite menu on unchanged data, with one
continuation. Arbitrary trainers, unforeseen menus, timing side channels
and comparisons with deployed central platforms remain outside the
evidence. Local tests assume trusted operations and durable state.

Written proofs, rational certificates and tests are distinct evidence.
Public experimental seeds and raw controls are not private packages.
No legal permissibility or lifelong inference ceiling is established.
Standard composition remains the fallback when no sharper complete-history
argument is available.

\section{Conclusion}
\label{sec:conclusion}
An existing export constrains what can be released next. Joint design
recovers utility beyond conditional composition in the tested finite
synthetic families. Compact exact certificates, cached public checks and
disjoint-block execution connect these mechanisms to committed releases.
They do not remove finite-state growth. The real-data test finds no
improvement over reuse; an information-value bound detects when further
private adaptation cannot achieve the prescribed gain over reuse. Practical gains for
richer models, unforeseen tasks and agency operation remain unestablished.

\bibliographystyle{mlsys2025}
\bibliography{references}
\clearpage
\appendix
\input{joint-commitment-proof.tex}
\input{model-history-proof.tex}
\input{count-extension-proof.tex}
\section{Executed release law and enforcement assumptions}
\label{app:finite-enforcement}
The finite protocol implements one continuation on an unchanged protected
state. Every task coordinate already exists before the old model is drawn;
newly acquired labels and changed records require another argument.

\paragraph{Commit the mechanism before protected input.}
The public specification fixes the old law $P_c(m)$, complete neighboring
count graph, task objectives, selector $T_c(s\mid m)$ and executable output
meanings. The actual channel menu is also verified and committed before
reservation or protected input. For each task $s$, exact checking establishes
\[
 Q_s(c,m,a)\ge0,\qquad \sum_a Q_s(c,m,a)=P_c(m),
\]
and $Q_s(c,m,a)\le B_s Q_s(c',m,a)$ for every declared neighbor.
Utility duals establish separate objective bounds; privacy needs only the
verified channel constraints. Objective binding alone is insufficient:
with a constant old output and zero objective, always-zero and always-one
channels are both optimal and zero-DP. Choosing between them using a private
bit discloses that bit. Consequently the optimizer and menu commitment use
only public or separately accounted information.

\paragraph{Law of the executed transcript.}
The selector has normalized nonnegative rows and satisfies
$T_c(s\mid m)\le\Xi T_{c'}(s\mid m)$ for each fixed old output.
The executor draws the old model, then the task, then the new model
conditional on the saved old output and unchanged confidential counts.
On a positive old-output branch,
\[
 \begin{aligned}
 &\Pr_c[M=m,S=s,A=a]\\
 &=P_c(m)T_c(s\mid m)\frac{Q_s(c,m,a)}{P_c(m)}\\
 &=T_c(s\mid m)Q_s(c,m,a).
 \end{aligned}
\]
Zero old mass forces zero joint mass and is never sampled. Multiplying
the two verified neighbor inequalities bounds the whole transcript by
$\Xi B_s\le\Xi\max_t B_t$. Summing gives pure DP with parameter
$\log(\Xi\max_t B_t)$, including the disclosed task. This does not
assert privacy of the conditional channel alone. Model descriptors and
copies are post-processing of this transcript. One committed continuation
per history prevents retries or additional recipients from obtaining
fresh uncharged samples.

\paragraph{Reservation and the common account.}
\label{app:central-gate-proof}
Before the old draw, reserve a conservative rational charge
$\epsilon_i\ge\log(\Xi\max_s B_s)$ for the complete history $i$.
Let $H$ contain distinct reserved or consumed histories and
$\kappa_i(u)$ their sufficient charges for unit $u$. Atomic admission
maintains
\[
 \sum_{i\in H}\kappa_i(u)\le E_u\qquad\text{for every }u.
\]
Consumption preserves this set; copying adds nothing, and revocation or
abandonment removes no charge. Induction over serialized transitions
establishes the invariant. Separate histories require valid conditional
composition under the preceding pooled history; independent approved
executions provide the supported route. Standard composition and group
privacy then justify the account~\citep{dwork2014foundations}.
A row-level history touching $r$ records of one person requires sufficient
charge $r\epsilon_i$ under the fixed-membership replacement contract.
Truthful linkage and complete contribution bounds are therefore essential.

\paragraph{Commitment, failure and current authorization.}
Reservation precedes execution; durable artifact commitment precedes
disclosure. Every possible output size is checked before any draw:
otherwise output-dependent rejection followed by retry could change the
certified law. Committed draws are immutable and reused after reopening;
failed uncommitted attempts reveal nothing. Recovery assumes faults do not
select on hidden random outcomes. Current source bindings, authority,
capacity and revocation are checked before mediated delivery. Successful
mathematical checks may be cached only for identical complete canonical
inputs, verifier identity and options, with byte equality after hashing.
The per-source digest omitted from shared mathematical templates remains
checked against the actual source. Incoming plans and menus are snapshotted;
a cache hit is never authorization.

\paragraph{Trust and preprocessing boundary.}
\label{app:central-audit}
\label{app:stage-zero}
The argument assumes trusted execution, entropy and durable authoritative
state, correctly encoded inputs, and mediation of every covered export.
Membership and pseudonymous linkage are fixed confidential context;
pseudonymization does not establish anonymity. Preprocessing must preserve
the declared neighboring relation or be included in the mechanism's privacy
argument. The prototype accepts aliases; it does not verify real identities
or implement a linkage vault. Policies may use public context and accounted
outputs. Unaccounted private decisions, diagnostics or timing cannot be
exposed under this argument. Internal records remain trusted control-plane
state. Local hashes do not prevent coherent ledger forgery or forks;
external export mediation and untrusted-host security are not established.

\input{richer-model-proof.tex}
\input{uncertain-workload-proof.tex}
\input{continuation-improvements-proof.tex}
\input{finite-realdata-proof.tex}
\subsection{Screening and reuse}
\label{app:screened-release}
This diagnostic was declared after the negative real-data study. It
does not replace that frozen evaluation. At minimum gain $1/200$, all
nine public problems fail the information-value screen: eight have
zero nonprivate value beyond reuse; TLC at $n=4$ permits only
$8.6893\times10^{-6}$. Enumerating all 256 old-output maps and an
attaining state-dependent oracle independently confirms these values.
The post-LP screen gives the same decisions.

One warm call followed by five timed calls gives median screen times
2.52--28.24\,ms on already constructed public problems, excluding public
fitting, domain construction and JSON loading. This avoids an LP, not
finite-state enumeration. The bound concerns the declared public prior
and output family, not unknown population accuracy.

Three separate $n=4$ workflows use actual calibration records, the four
binary-feature heads and a reuse certificate constructed directly from
the public problem, with LP calls disabled. Registration and exact
validation precede the old draw. A prospective charge and cap of
$78/125=.624$ cover old-law odds $97/52$; an odds-$3$ plan is rejected.
Committed model bytes survive cold-cache ledger reopening, retries and
delivery to another recipient. Another charged history is refused at
the exhausted cap. Revocation blocks subsequent delivery but retains
the charge. An identical-channel control retains its earlier $1.1$
reservation after revocation: tighter prospective accounting neither
improves an unchanged channel's statistical privacy nor refunds history.
These local checks establish no private-learning gain or agency deployment.

\end{document}
